% Options for packages loaded elsewhere
\PassOptionsToPackage{unicode}{hyperref}
\PassOptionsToPackage{hyphens}{url}
%
\documentclass[
  11pt,
]{article}
\usepackage{lmodern}
\usepackage{amssymb,amsmath}
\usepackage{ifxetex,ifluatex}
\ifnum 0\ifxetex 1\fi\ifluatex 1\fi=0 % if pdftex
  \usepackage[T1]{fontenc}
  \usepackage[utf8]{inputenc}
  \usepackage{textcomp} % provide euro and other symbols
\else % if luatex or xetex
  \usepackage{unicode-math}
  \defaultfontfeatures{Scale=MatchLowercase}
  \defaultfontfeatures[\rmfamily]{Ligatures=TeX,Scale=1}
\fi
% Use upquote if available, for straight quotes in verbatim environments
\IfFileExists{upquote.sty}{\usepackage{upquote}}{}
\IfFileExists{microtype.sty}{% use microtype if available
  \usepackage[]{microtype}
  \UseMicrotypeSet[protrusion]{basicmath} % disable protrusion for tt fonts
}{}
\makeatletter
\@ifundefined{KOMAClassName}{% if non-KOMA class
  \IfFileExists{parskip.sty}{%
    \usepackage{parskip}
  }{% else
    \setlength{\parindent}{0pt}
    \setlength{\parskip}{6pt plus 2pt minus 1pt}}
}{% if KOMA class
  \KOMAoptions{parskip=half}}
\makeatother
\usepackage{xcolor}
\IfFileExists{xurl.sty}{\usepackage{xurl}}{} % add URL line breaks if available
\IfFileExists{bookmark.sty}{\usepackage{bookmark}}{\usepackage{hyperref}}
\hypersetup{
  pdftitle={Polarized closure of the integral Einstein obstruction},
  hidelinks,
  pdfcreator={LaTeX via pandoc}}
\urlstyle{same} % disable monospaced font for URLs
\usepackage[margin=1in]{geometry}
\setlength{\emergencystretch}{3em} % prevent overfull lines
\providecommand{\tightlist}{%
  \setlength{\itemsep}{0pt}\setlength{\parskip}{0pt}}
\setcounter{secnumdepth}{5}
\usepackage{mathrsfs}
\setlength{\emergencystretch}{3em}

\title{Polarized closure of the integral Einstein obstruction}
\usepackage{etoolbox}
\makeatletter
\providecommand{\subtitle}[1]{% add subtitle to \maketitle
  \apptocmd{\@title}{\par {\large #1 \par}}{}{}
}
\makeatother
\subtitle{Gauge input, uniform constants, and all Einstein signs}
\author{}
\date{22 September 2026}

\begin{document}
\maketitle

\hypertarget{statement-and-analytic-dependency}{%
\section{Statement and analytic
dependency}\label{statement-and-analytic-dependency}}

Let \(X\) be a closed connected oriented smooth four-manifold, let

\[
T_X=2\chi(X)+3\tau(X)>0,
\qquad M=X\#k\overline{\mathbb{CP}}^{\,2}
                 \#\ell(S^1\times S^3),\qquad D=k+4\ell.
\]

Assume that a Spin\(^c\) structure with real first Chern class \(c\) has
a nonzero ordinary Seiberg--Witten invariant, with its insertion
retained if the expected dimension is positive. When \(b_+(X)=1\),
nonvanishing is in the small-perturbation chamber, not an arbitrarily
chosen chamber. Write \(A=c^2\ge T_X>0\), and define

\[
\mathcal E(g)=\inf_{\lambda\ne0}
 \frac{\int_M|\operatorname{Ric}_g-\lambda g|_g^2\,d\mu_g}
 {\lambda^2\operatorname{Vol}_g(M)}.
\]

\textbf{Stability theorem, relative to the integral calibration lemma
below.} There is a universal constant \(\varepsilon>0\) such that every
smooth metric on \(M\) satisfies both

\[
\mathcal E(g)+\left(\frac{T_X-D}{A}-\frac23\right)_+
 \ge\varepsilon,
\qquad
\mathcal E(g)+\left(\frac13-\frac D{T_X}\right)_+
 \ge\varepsilon.                                      \tag{1}
\]

In particular, if \(X\) is any minimal complex surface of general type,
including \(b_+(X)=1\), then \(A=T_X=c_1^2(X)>0\), and

\[
\boxed{\mathcal E(g)+
 \left(\frac13-\frac{k+4\ell}{c_1^2(X)}\right)_+
 \ge\varepsilon.}                                    \tag{2}
\]

This note proves the gauge-to-obstruction implication, including its
constants and the zero Einstein constant. Its nonclassical analytic
input is stated precisely in Section 3. The full proof of that lemma
remains the integral calibration manuscript; restating it here is not an
independent verification of its Hessian and weak-Bianchi arguments.
Neither a universal supply of special monopole classes on arbitrary
four-manifolds nor positivity of a universal special-class invariant is
assumed.

\hypertarget{the-actual-polarized-class-and-every-auxiliary-weight}{%
\section{The actual polarized class and every auxiliary
weight}\label{the-actual-polarized-class-and-every-auxiliary-weight}}

\hypertarget{sign-selection-and-the-chamber}{%
\subsection{Sign selection and the
chamber}\label{sign-selection-and-the-chamber}}

Fix \(g\), and let \(H=H_g^+\) be its positive harmonic subspace. Pull
\(c\) back to \(M\). Choose exceptional generators \(E_{j,H}\),
\textbf{with signs chosen at this \(H\)}, so that

\[
E_{i,H}\cdot E_{j,H}=-\delta_{ij},\qquad
c\cdot E_{j,H}=0,\qquad c_H^+\cdot E_{j,H}\le0.
\]

Polarized Seiberg--Witten gluing supplies an actual Spin\(^c\) structure
with

\[
a_H=c-\sum_j E_{j,H}                                  \tag{3}
\]

and nonzero invariant at \(H\). Thus its unperturbed equations are
solvable for every \(H\)-adapted metric. The classical input is the
construction in LeBrun's Proposition 3.2 {[}L01{]}, not a consequence
merely of the following intersection-form calculation:

\[
\begin{split}
(a_H^+)^2
 &=(c_H^+)^2-2\sum_j c_H^+\cdot E_{j,H}
                  +\left|\sum_j E_{j,H}^+\right|^2
 \ge(c_H^+)^2\ge A,\\
a_H^2&=A-k\ge T_X-k-4\ell=T_M.
\end{split}                                           \tag{4}
\]

For \(b_+=1\), choose the positive generator \(\omega_H\) so that
\(c\cdot\omega_H>0\). This is possible because \(c^2>0\). The sign
selection implies

\[
a_H\cdot\omega_H=c\cdot\omega_H-\sum_j E_{j,H}\cdot\omega_H>0.
\]

The corresponding base projection of \(\omega_H\) has positive square.
Hence the small-perturbation chamber agrees with the one in the
chamber-compatible blow-up formula. Handle additions do not change the
intersection form. If the base invariant uses an insertion \(z\), retain
it and add the \(\ell\) degree-one loop insertions. Indeed,

\[
\frac{a_H^2-T_M}{4}=\frac{A-T_X}{4}+\ell.
\]

The surgery and insertion statements used here are classical gauge
theory {[}OS{]}; this dimension equality alone would not prove
nonvanishing.

For a minimal general-type complex surface, the canonical Spin\(^c\)
structure has \(c=-K_X\), \(K_X^2>0\), and \(K_X\cdot[\omega]>0\) for a
Kähler form. When \(b_+=1\), put \(K_X\) in the forward cone. Its
pairing with every unit forward positive generator is at least
\(\sqrt{K_X^2}>0\). The zero-perturbation discriminant for \(-K_X\) is
therefore negative, identifying the small chamber with the canonical
symplectic chamber where nonvanishing holds {[}LL{]}. This is why
\(b_+=1\) is included. A symplectic structure without this chamber
verification would not suffice.

\hypertarget{conformal-covariance-with-the-bundle-identification-explicit}{%
\subsection{Conformal covariance, with the bundle identification
explicit}\label{conformal-covariance-with-the-bundle-identification-explicit}}

Let \(\mathfrak q_g\) denote the quadratic spinor-to-self-dual-two-form
map, normalized by \(|\mathfrak q_g(\Phi)|^2=|\Phi|^4/8\). For any
smooth \(f>0\), set

\[
\widehat g=f^{-2}g=e^{2u}g,\qquad u=-\log f.
\]

The Hodge star on two-forms is unchanged, so
\(H_{\widehat g}^+=H_g^+=H\). Use the \textbf{same} Spin\(^c\) structure
from (3) to solve the unperturbed equations for \(\widehat g\):

\[
D_{\mathcal A}^{\widehat g}\widehat\Phi=0,\qquad
-iF_{\mathcal A}^+=\mathfrak q_{\widehat g}(\widehat\Phi).
\]

Let \(\beta:S_g\to S_{\widehat g}\) be the fibrewise unitary spinor
identification. The conformal identities are

\[
D_{\mathcal A}^{\widehat g}(e^{-3u/2}\beta\Phi)
 =e^{-5u/2}\beta D_{\mathcal A}^{g}\Phi,
\qquad
\mathfrak q_{\widehat g}(\beta v)=e^{2u}\mathfrak q_g(v),
\]

where the second equality compares covariant two-forms. With
\(\Phi=f^{-3/2}\beta^{-1}\widehat\Phi\), they give

\[
D_{\mathcal A}^{g}\Phi=0,\qquad
\mathfrak q_g(\Phi)=f^{-1}\mathfrak q_{\widehat g}(\widehat\Phi),
\qquad -iF_{\mathcal A}^+=f\mathfrak q_g(\Phi).           \tag{5}
\]

This is the conformal computation underlying {[}L06, Lemma 3.7{]},
applied only to the one conformal class needed here. The exact order is

\[
\forall g\ \exists\mathfrak s_g\ \forall f\in C^\infty(M,(0,\infty))
\ \exists(\mathcal A_f,\Phi_f)\text{ solving (5)}.       \tag{6}
\]

No continuity of \(g\mapsto\mathfrak s_g\) is required. The class may
change between metrics; it is held fixed while \(f\) varies. Ordinary
solvability just at \(g\) would not establish (6).

\hypertarget{the-analytic-input-and-the-uniform-constant}{%
\section{The analytic input and the uniform
constant}\label{the-analytic-input-and-the-uniform-constant}}

All Ricci norms below are tensor Hilbert--Schmidt norms. The norm on
\(W^+\) is its Hilbert--Schmidt norm as an endomorphism of
\(\Lambda^+\). Put \(V=\operatorname{Vol}(g)\), \(T=T_M\), and

\[
\begin{gathered}
\mathscr O=\operatorname{Orb}\operatorname{diag}(-1/6,1/12,1/12),
F=\operatorname{dist}(W^+,\mathscr O)^2,\qquad
q=-s+\sqrt6|W^+|,\qquad p=(a_H^+)^2,\\
\Delta=\int q_+^2-72\pi^2p,\qquad
\kappa=\frac1V\left(\int(|\operatorname{Ric}+g|^2+F)+\Delta\right),
\qquad B=\frac{(-a_H^2)_+}{V}.
\end{gathered}
\]

\textbf{Integral calibration lemma (the nonclassical input).} For an
actual structure satisfying (6) and \(p>0\), the signed spectral
monopole inequality gives \(\Delta\ge0\), and a universal
\(c_{\rm cal}>0\) satisfies

\[
                    (1+B)\kappa\ge c_{\rm cal}.         \tag{7}
\]

The proof in the integral calibration manuscript constructs bounded
weighted spinors, obtains a covariant Hessian estimate using integral
Ricci and connection curvature, forms a smooth polynomial Weyl tensor,
and tests weak Bianchi against its divergence. Its hypotheses involve
the selected structure at the metric under consideration, not an
all-metric invariant. This note uses (7) explicitly; it does not replace
that proof by an appeal to a previous round's conclusion.

We next give the complete constant transfer. Put

\[
r=V^{-1}\int|\operatorname{Ric}+g|^2,\qquad
d=(T/A-2/3)_+.
\]

The algebraic comparison detailed in the appendix yields

\[
\kappa\le\frac{245}{2}r+972\pi^2\frac AVd.              \tag{8}
\]

The scalar monopole bound gives

\[
32\pi^2A\le\int s_-^2\le\int s^2\le32V+8rV,
\qquad A/V\le(4+r)/(4\pi^2).
\]

Consequently \(r\le1\) implies \(\kappa\le1215(r+d)\). For \(r\le1/4\),
Chern--Weil and the reverse triangle inequality give

\[
\begin{split}
4\pi^2T/V
&=V^{-1}\int\left(s^2/24+2|W^+|^2
                    -|\operatorname{Ric}^{\circ}|^2/2\right)\\
&\ge(4-2\sqrt r)^2/24-r/2\ge1/4.
\end{split}                                           \tag{9}
\]

Since \(a_H^2\ge T>0\), we have \(B=0\) in this branch, even when the
family allows arbitrarily many blow-ups. Equations (7)--(9) imply

\[
r+d\ge\varepsilon,\qquad
\varepsilon=\min\{1/4,c_{\rm cal}/1215\}>0.             \tag{10}
\]

For \(r>1/4\), (10) is immediate. This proves uniformity with respect to
the manifold, the metric, the number of exceptional classes, and their
selection. Smooth approximations of the auxiliary weight are taken with
\(g,a_H\) fixed before their approximation error tends to zero; no
uniform rate of smoothing across metrics is used.

For the positive normalization, set
\(r_+=V^{-1}\int|\operatorname{Ric}-g|^2\). On the scalar-negative set,
\(s_-^2\le(s-4)^2\), whence

\[
A/V\le r_+/(8\pi^2).
\]

If \(r_+\le1/4\) and \(d\le1/4\), the analogue of (9) gives
\(4\pi^2T/V\ge1/4\), whereas \(T\le A(2/3+d)\) gives

\[
4\pi^2T/V\le(r_+/2)(2/3+d)\le11/96<1/4.
\]

Thus \(r_++d>1/4\). For any \(\lambda\ne0\), apply the appropriate
normalization to \(\widetilde g=|\lambda|g\). Constant rescaling
preserves \(H\), (4), and the weighted input. In dimension four,

\[
\frac1{V_{\widetilde g}}\int
 |\operatorname{Ric}_{\widetilde g}-\operatorname{sign}(\lambda)
             \widetilde g|_{\widetilde g}^2
=\frac{\int|\operatorname{Ric}_g-\lambda g|_g^2}{\lambda^2V_g}.
\]

Taking the infimum proves the first inequality in (1). Apply the same
argument with the smaller positive projection budget \(T_X\le A\) to
obtain the second. No comparison of ratios with a potentially negative
numerator is necessary. Equation (2) follows by \(A=T_X\).

\hypertarget{scalar-flat-exclusion-and-all-einstein-constants}{%
\section{Scalar-flat exclusion and all Einstein
constants}\label{scalar-flat-exclusion-and-all-einstein-constants}}

With \(f=1\), the spinor Weitzenböck formula, integrated on the closed
manifold, gives

\[
0=\int(4|\nabla_{\mathcal A}\Phi|^2+s|\Phi|^2+|\Phi|^4).
\]

Writing \(\rho=|\Phi|^2\), Cauchy--Schwarz and the curvature equation
yield

\[
\int s_-^2\ge\int\rho^2
=8\int|F_{\mathcal A}^+|^2
\ge32\pi^2(a_H^+)^2\ge32\pi^2A>0.                     \tag{11}
\]

The spinor cannot vanish identically because \(p>0\). In particular,
there are no scalar-flat metrics in this family, for any \(k,\ell\).

For \(\operatorname{Ric}_g=\lambda_0g\) with \(\lambda_0\ne0\), the
definition gives \(\mathcal E(g)=0\). For a Ricci-flat metric it instead
gives \(\mathcal E(g)=4\), since \(|\lambda g|_g^2=4\lambda^2\). The
Ricci-flat case is therefore excluded by (11), not by assigning it zero
defect.

Choose \(\gamma=\varepsilon/2\). From (1) and (11),

\[
\boxed{D\ge(1/3-\gamma)T_X
\quad\Longrightarrow\quad
M\text{ admits no Einstein metric of any sign}.}        \tag{12}
\]

The sharper necessary inequality for a nonzero Einstein constant is
\(D/T_X\le1/3-\varepsilon\). The closed threshold (12) uses a smaller
constant intentionally; equality at \(1/3-\varepsilon\) is not excluded
by (1) alone.

For general-type \(X\) and \(\ell=0\), (12) contains every case of
LeBrun's Corollary 3.4. In its closed band, every metric also satisfies

\[
\mathcal E(g)\ge\varepsilon/2,\qquad
\delta_\sigma(g)\ge\sqrt{\varepsilon/8},\quad
\delta_\sigma(g)=\|g^{-1}\operatorname{Ric}_g-\sigma I\|_{\infty,\rm op},
\quad \sigma=\pm1,
\]

because \(\mathcal E(g)\le4\delta_\sigma(g)^2\).

\hypertarget{theorem-3.3-the-positive-budget-and-the-zero-corner}{%
\section{Theorem 3.3: the positive budget and the zero
corner}\label{theorem-3.3-the-positive-budget-and-the-zero-corner}}

The positive-budget connected-sum setting of Theorem 3.3 is covered by
(1) and (12), including its handle term and the specified chamber. No
new gauge theorem is required beyond its polarized input.

For nonpositive budgets, use Chern--Weil directly. If \(T_M<0\), the
identity in (9) excludes Einstein metrics. There is also an elementary
integral bound that does not require any monopole input. Set

\[
R_2=\int|\operatorname{Ric}|^2,\quad S=\int s,\quad
L=\int(s-S/V)^2,\quad W_2=\int|W^+|^2.
\]

When \(R_2>0\), optimization in \(1/\lambda\) gives

\[
\mathcal E=4-\frac{S^2}{VR_2}
=1+\frac{12W_2-24\pi^2T_M+L}{R_2}.                     \tag{13}
\]

Thus \(T_M\le0\) implies \(\mathcal E\ge1\), with the case \(R_2=0\)
handled by \(\mathcal E=4\). This is a classical algebraic consequence,
not a new calibration theorem.

If \(T_X<0\), or \(T_X=0<D\), then \(T_M=T_X-D<0\), and the exact
nonexistence conclusion follows. However, \textbf{\(T_X=D=0\) is
different}: an Einstein metric is then forced to be Ricci-flat and
anti-self-dual, not impossible. K3 with its complex orientation is an
actual example. It has nonzero Seiberg--Witten invariant and admits
Ricci-flat metrics.

The inspected author-hosted PDF and arXiv v1 of {[}L01{]} are the March
2000 preprint, not the journal typesetting. Their unrestricted wording
of Theorem 3.3 must not be imported across this zero corner. The
journal's publication metadata alone does not establish how that corner
was worded in its final text. Our precise claim is the quantitative
extension for \(T_X>0\), together with the classical \(T_M<0\)
exclusion. This qualification does not narrow Corollary 3.4: there
\(T_X=c_1^2(X)>0\).

\hypertarget{what-is-strengthened}{%
\section{What is strengthened}\label{what-is-strengthened}}

The additional threshold interval contains actual members of the
general-type family. Let

\[
X_m=\Sigma_2\times\Sigma_m,\qquad A_m=8(m-1),\qquad
k_m=\lceil A_m/3\rceil-1,\qquad m\ge2.
\]

These are minimal general-type bases and

\[
0<1/3-k_m/A_m\le1/A_m\longrightarrow0.
\]

For all sufficiently large \(m\), their blow-ups satisfy (12), although
\(k_m<A_m/3\). Their Euler characteristics tend to infinity. Thus the
strengthening is not merely an enlargement of unattained real parameter
values. This does not assert novelty over every other obstruction in the
literature, nor produce an explicit first \(m\) without an explicit
calibration constant.

The correct implication is \textbf{integral stability plus scalar-flat
exclusion implies the classical exact obstruction}. The exact
obstruction alone would not imply a uniform integral gap. LeBrun's
absolute Ricci-energy infima and minimal-volume results are different
conclusions, not consequences established by this note.

\hypertarget{appendix-algebra-behind-the-constant-transfer}{%
\section{Appendix: algebra behind the constant
transfer}\label{appendix-algebra-behind-the-constant-transfer}}

Write \(e=s+4\), \(\alpha=\sqrt{2/3}|W^+|\), \(w=\lambda_{\min}(W^+)\),
\(t=s+24\alpha\), and

\[
b=2s/3+2w,\quad
G=(2s/3-2\alpha)^2-b_-^2\ge0,\quad
P=\int|\operatorname{Ric}^{\circ}|^2.
\]

The signed spectral monopole bound is \(\int b_-^2\ge32\pi^2p\). Here
\(v_+=\max(v,0)\) and \(v_-=\max(-v,0)\) for a real number \(v\). The
pointwise identity

\[
t^2+18G=216(s^2/24+2|W^+|^2)-18b_-^2
\]

therefore gives

\[
\int(t^2+18G)\le108P+864\pi^2Ad.
\]

Also \(F\le e^2+t^2+G\). One direct verification sets
\(z=w+\alpha\in[0,\alpha]\), so

\[
F=\tfrac32(\alpha-1/6)^2+z/2,\qquad
\alpha=(t-e+4)/24.
\]

For \(s\le-3\), \(b<0\) and \(G\ge8z\), giving
\(F\le(t^2+e^2)/192+G/16\). For \(s>-3\), \(e>1\), and \(z\le\alpha\),
\(4-e\le3e^2\), \(t\le(t^2+1)/2\) give \(F\le(t^2+5e^2)/64\). Hence

\[
\int(|\operatorname{Ric}+g|^2+F)
 \le109rV+864\pi^2Ad.
\]

Finally, \(27(s^2/24+2|W^+|^2)-q^2=t^2/8\), so \(q_+^2\le q^2\) and
\(p\ge A\) yield

\[
\Delta\le(27/2)rV+108\pi^2Ad.
\]

Adding proves (8). These estimates involve only the selected class's
projection bound and its square; neither a finite class set nor a
continuous selection is hidden in the constants.

\hypertarget{references-and-proof-locations}{%
\section{References and proof
locations}\label{references-and-proof-locations}}

\begin{itemize}
\tightlist
\item
  {[}L01{]} C. LeBrun, \emph{Ricci Curvature, Minimal Volumes, and
  Seiberg--Witten Theory}, Invent. Math. \textbf{145} (2001), 279--316.
  \href{https://www.math.stonybrook.edu/~claude/min.pdf}{Author-hosted
  preprint}, Definition 1, Proposition 3.2, Theorem 3.3, Corollary 3.4;
  \href{https://doi.org/10.1007/s002220100148}{publication record}.
\item
  {[}L06{]} C. LeBrun, \emph{Four-manifolds, curvature bounds, and
  convex geometry}, \href{https://arxiv.org/abs/math/0611450}{preprint},
  Lemma 3.7.
\item
  {[}LL{]} T.-J. Li and A.-K. Liu, \emph{Uniqueness of symplectic
  canonical class, surface cone and symplectic cone of 4-manifolds with
  \(b^+=1\)},
  \href{https://arxiv.org/abs/math/0012048}{arXiv:math/0012048}, Section
  2.1.
\item
  {[}OS{]} P. Ozsváth and Z. Szabó, \emph{Higher type adjunction
  inequalities in Seiberg--Witten theory},
  \href{https://arxiv.org/abs/math/0005268}{arXiv:math/0005268},
  Proposition 2.2, Remark 2.3, Theorem 3.2.
\end{itemize}

The package's analytic proof of (7) is
\href{round4_rigidity.md}{Integral calibration}, Sections 1--6; the
exact spinor and weighted Hodge identities are in
\href{round3_monopole_input.md}{Monopole input}, Sections 6--8. The
earlier \href{round13_chamber.md}{polarized proof} and
\href{round21_chambers.md}{chamber analysis} remain intact. The present
note consolidates their interfaces; it is not a new independent main
theorem or an audit of every analytic identity in those manuscripts.

\end{document}
